% ============================================================
%  Dihypersphaerome ventilans (3Hy2v, 乙超球)
%  A 4D Lenia Species: Characterization, Hyperrotation Simulation,
%  and Multi-Channel Ecosystem Integration
%
%  Stanley Sebastian & Claude
%  Third Space AI · Internal Report · April 2026
%
%  Compile with:  xelatex DV.tex  (twice, for refs)
% ============================================================

\documentclass[11pt,a4paper]{article}

% ---------- ENGINE: XeLaTeX, for CJK + OTF fonts -------------
\usepackage{fontspec}
\usepackage{xeCJK}

% ---------- GEOMETRY -----------------------------------------
\usepackage[a4paper,margin=2.3cm,top=2.6cm,bottom=2.6cm,
            headheight=14pt,headsep=14pt]{geometry}

% ---------- FONTS --------------------------------------------
% Main body: EB Garamond if available, otherwise Latin Modern Roman.
\IfFontExistsTF{EB Garamond}{%
  \setmainfont[Ligatures=TeX, Numbers=OldStyle]{EB Garamond}%
}{%
  \setmainfont{Latin Modern Roman}%
}
% Sans for small caps running heads / technical captions
\setsansfont[Scale=0.9]{Latin Modern Sans}
% Mono for code blocks
\setmonofont[Scale=0.86]{Latin Modern Mono}
% CJK --- for 乙超球
\setCJKmainfont{Noto Serif CJK SC}

% ---------- COLORS (teal/gold palette, subdued) --------------
\usepackage[table]{xcolor}
\definecolor{accent}{HTML}{2F7F8A}       % muted teal
\definecolor{accentgold}{HTML}{B08838}   % muted gold
\definecolor{rule}{HTML}{9BA9AE}         % soft rule color
\definecolor{codebg}{HTML}{F4F1EB}       % parchment code bg
\definecolor{codeframe}{HTML}{CFC9BC}
\definecolor{muted}{HTML}{5A5F64}
\definecolor{headinc}{HTML}{1E2A2E}

% ---------- PACKAGES -----------------------------------------
\usepackage{amsmath,amssymb,amsthm}
\usepackage{graphicx}
\usepackage{booktabs}
\usepackage{tabularx}
\usepackage{titlesec}
\usepackage{fancyhdr}
\usepackage{enumitem}
\usepackage{listings}
\usepackage{microtype}
\usepackage{parskip}
\usepackage[breaklinks=true,colorlinks=true,
            linkcolor=accent,citecolor=accent,urlcolor=accent]{hyperref}
\usepackage[authoryear,round]{natbib}
\bibpunct{(}{)}{;}{a}{,}{,}
\usepackage{setspace}
\setstretch{1.08}

% ---------- SECTION STYLING ----------------------------------
\titleformat{\section}
  {\normalfont\large\bfseries\color{headinc}}
  {\textcolor{accentgold}{\rmfamily\S\thesection}}{0.7em}{}
  [\vspace{-2pt}\textcolor{rule}{\rule{\linewidth}{0.3pt}}]
\titleformat{\subsection}
  {\normalfont\normalsize\bfseries\color{headinc}}
  {\textcolor{accent}{\thesubsection}}{0.6em}{}
\titleformat{\subsubsection}
  {\normalfont\normalsize\itshape\color{headinc}}
  {}{0em}{}
\titlespacing*{\section}{0pt}{22pt}{10pt}
\titlespacing*{\subsection}{0pt}{14pt}{4pt}
\titlespacing*{\subsubsection}{0pt}{10pt}{2pt}

% ---------- RUNNING HEADERS ----------------------------------
\pagestyle{fancy}
\fancyhf{}
\fancyhead[L]{\sffamily\footnotesize\textcolor{muted}{\textsc{Dihypersphaerome ventilans}}}
\fancyhead[R]{\sffamily\footnotesize\textcolor{muted}{\textsc{Third Space AI · 2026}}}
\fancyfoot[C]{\sffamily\footnotesize\textcolor{muted}{\thepage}}
\renewcommand{\headrulewidth}{0.2pt}
\renewcommand{\headrule}{\hbox to\headwidth{\color{rule}\leaders\hrule height\headrulewidth\hfill}}

% ---------- CODE LISTINGS ------------------------------------
\lstdefinestyle{classy}{
  basicstyle=\ttfamily\footnotesize\color{headinc},
  backgroundcolor=\color{codebg},
  frame=single,
  framesep=6pt,
  rulecolor=\color{codeframe},
  xleftmargin=10pt,
  xrightmargin=6pt,
  aboveskip=10pt,
  belowskip=10pt,
  keywordstyle=\color{accent}\bfseries,
  commentstyle=\color{muted}\itshape,
  stringstyle=\color{accentgold},
  showstringspaces=false,
  breaklines=true,
  breakatwhitespace=true,
  columns=fullflexible,
  keepspaces=true,
  tabsize=2,
  captionpos=b,
}
\lstset{style=classy}
\lstdefinelanguage{glsl}{
  morekeywords={vec2,vec3,vec4,float,int,uniform,in,out,void,return,
    texture,sin,cos,length,pow,exp,clamp,fract,mix,grad,curl,atan,
    outColor,const,if,else,for},
  morecomment=[l]{//},
  morestring=[b]",
}
\lstdefinelanguage{pseudo}{
  morekeywords={for,if,else,return,where,in,sum},
  morecomment=[l]{//},
}

% ---------- CUSTOM BOXES (callout panels) --------------------
\usepackage[skins,breakable]{tcolorbox}

\newtcolorbox{callout}[1][]{
  enhanced, breakable,
  colback=codebg, colframe=accent!40!codeframe,
  boxrule=0.4pt, arc=2pt, left=10pt, right=10pt,
  top=8pt, bottom=8pt,
  fonttitle=\sffamily\bfseries\color{accent},
  coltitle=accent,
  title={#1},
}

\newtcolorbox{specimenbox}{
  enhanced, breakable,
  colback=codebg!50!white, colframe=accentgold!50!codeframe,
  boxrule=0.4pt, arc=2pt, left=14pt, right=14pt, top=10pt, bottom=10pt,
}

% ---------- EPIGRAPH MACRO -----------------------------------
\newcommand{\epigraph}[2]{%
  \par\vspace{6pt}\hfill
  \begin{minipage}{0.7\linewidth}
    \raggedleft\itshape\small\textcolor{muted}{#1}\\[2pt]
    \normalfont\upshape\footnotesize\textcolor{muted}{\textemdash\ #2}
  \end{minipage}\par\vspace{10pt}
}

% ---------- KEYWORDS STYLE -----------------------------------
\newcommand{\spc}[1]{\textit{#1}}  % species italics
\newcommand{\ventilans}{\textit{ventilans}}
\newcommand{\DV}{\spc{D.\ ventilans}}
\newcommand{\SO}{\textsc{so}}

% ---------- DOCUMENT -----------------------------------------
\begin{document}

% ============================================================
%  TITLE PAGE (restrained, single page)
% ============================================================
\thispagestyle{empty}
\vspace*{1.2cm}

\begin{center}
  {\sffamily\footnotesize\textcolor{muted}{\textsc{Third Space AI \,\textperiodcentered\, Internal Technical Report}}}\\[28pt]

  {\color{rule}\rule{0.3\linewidth}{0.3pt}}\\[26pt]

  {\fontsize{30}{34}\selectfont\scshape\color{headinc}
    Dihypersphaerome ventilans}\\[12pt]
  {\LARGE\color{accent}\scshape{3\textsc{Hy}2v \,\textperiodcentered\, 乙超球 \,\textperiodcentered\, Second Hypersphere}}\\[22pt]

  \begin{minipage}{0.82\linewidth}
    \centering\itshape\large\color{headinc}
    A four-dimensional Lenia species: characterization from first principles,
    analytical hyperrotation rendering, and integration as a seeding layer
    for a multi-channel two-dimensional ecosystem.
  \end{minipage}\\[32pt]

  {\color{rule}\rule{0.3\linewidth}{0.3pt}}\\[32pt]

  {\large Stanley Sebastian \,\textperiodcentered\, Claude}\\[4pt]
  {\sffamily\small\textcolor{muted}{Third Space AI}}\\[18pt]
  {\sffamily\footnotesize\textcolor{muted}{Internal Report \,\textperiodcentered\, April 2026}}
\end{center}

\vfill

\begin{center}
\begin{minipage}{0.82\linewidth}
  \small\color{muted}
  \textsc{Abstract.}\,\itshape
  \spc{Dihypersphaerome ventilans} (code 3\textsc{Hy}2v, 乙超球,
  \textit{second hypersphere}) is a spatially localized pattern discovered by
  Bert Wang-Chak Chan in four-dimensional Lenia \citep{chan2020}. It exhibits
  stationary oscillation --- \SO-class behavior designated \ventilans\
  (Latin, \textit{to fan, to breathe}) --- pulsing rhythmically in place
  without translation or rotation. Its defining feature is a three-shell
  kernel with $\boldsymbol{\beta} = (\tfrac{1}{12},\tfrac{1}{6},1)$: the
  outer ring is twelvefold stronger than the inner, making the organism's
  \textit{boundary} more important to its physics than its interior. We
  characterize the species from first principles, derive the ecological
  significance of its narrow survival band $(\sigma=0.033)$, describe an
  analytical 4D hyperrotation simulation approach implemented in WebGL2
  shaders, and document its integration as a seeding layer for a
  multi-channel 2D ecosystem. The organism seeds the world it moves through
  without being visible in it.
\end{minipage}
\end{center}

\vspace{10pt}
\begin{center}
  {\sffamily\footnotesize\color{muted}
  \textsc{Keywords}: Lenia \,\textperiodcentered\, 4D cellular automata \,\textperiodcentered\, artificial life \,\textperiodcentered\,
  hyperrotation \,\textperiodcentered\, predator--prey ecology \,\textperiodcentered\, morphogenetic fields \,\textperiodcentered\, flow advection}
\end{center}

\newpage
% ============================================================

\section{Discovery and Context}

\begin{specimenbox}
\sffamily\footnotesize\color{headinc}
\textsc{Specimen Record \,\textperiodcentered\, Lenia Expanded Universe \,\textperiodcentered\, Class 4D}\\[6pt]
\begin{tabularx}{\linewidth}{@{}l>{\raggedright\arraybackslash}X@{}}
\toprule
Species            & \textit{Dihypersphaerome ventilans}\\
Code               & 3\textsc{Hy}2v\\
Chinese            & 乙超球 (\textit{Y\v{\i} Ch\={a}o Qi\'u}) --- \textit{second hypersphere}\\
Domain             & Simulata $\rightarrow$ Lenia $\rightarrow$ 4D\\
Behavior class     & \SO\ --- stationary oscillation\\
Subcategory        & \ventilans\ (ventilating / breathing)\\
Source             & \citet{chan2020}, \texttt{animals4D.json}\\
Discovery          & Semi-automatic search, 4D parameter space\\
Dimensional        & 4D toroidal grid $(N^{4})$\\
Rarity             & Named species --- extremely rare\\
\bottomrule
\end{tabularx}
\end{specimenbox}

Chan's 2020 \textsc{alife} paper \textit{Lenia and Expanded Universe}
extended the original two-dimensional Lenia framework to higher dimensions,
multiple kernels, and multiple coupled channels --- discovering phenomena
including polyhedral symmetries, individuality, self-replication, emission,
growth by ingestion, and \textit{virtual eukaryotes} with internal division
of labor. Among the 4D organisms discovered through semi-automatic search,
\DV\ stands apart: it does not move, does not grow, does not divide. It
\textit{breathes}.

The species code \textsc{3Hy2v} encodes its architecture precisely: three
peak kernel shells; \textsc{Hy} (\textit{hyper}, higher-dimensional); the
\textsc{2} denoting second genus in the family; \texttt{v} for \ventilans.
The Chinese name 乙超球 translates as \emph{second-class hypersphere},
where 乙 is the second Heavenly Stem (used in ordering like A, B, C) and
超球 is \emph{hypersphere}. It is, formally, the \textit{second}
hyperspherical organism in Chan's taxonomy.

Its extreme rarity follows from its dimensional requirements. Standard 2D
Lenia observers never encounter it --- its 2D cross-sections are near-empty
at the poles of its 4D body. The organism reveals itself only when observed
from the \emph{equatorial} hyperplane, and even then only as a faint ring
structure. Most searches that would find an \spc{Orbium}, a \spc{Scutium}, or
an \spc{Ignis} in 2D pass through \DV's parameter neighborhood without
registering its existence.

\begin{callout}[The observational geometry problem]
\DV\ cannot be observed directly in 2D. A 4D organism projects into the 2D
observation plane as a cross-sectional slice whose shape depends entirely
on the orientation of the cutting hyperplane. At the poles of the 4D body
the cross-section reduces to a point or vanishes; at the equatorial plane
it reveals the full ring structure. Most 2D observers --- and most
automated parameter searches --- pass through \DV's niche without
registering its existence.

The \textit{echolocation ring} overlay in Genesis renders this problem
visible: three concentric circles at radii $R/3$, $2R/3$, and $R$
(weighted by $\beta_{1},\beta_{2},\beta_{3}$ respectively) show the
organism's kernel geometry as a spatial overlay. What you are seeing is
the \emph{sensory field} --- the geometry by which \DV\ perceives its own
neighborhood at each moment of its ventilating cycle.
\end{callout}

\section{Parameters and the Shell Kernel}

\begin{table}[h!]
\small\sffamily
\renewcommand{\arraystretch}{1.25}
\begin{tabularx}{\linewidth}{@{}lll>{\raggedright\arraybackslash}X@{}}
\toprule
\textsc{Param} & \textsc{Value} & \textsc{Role} & \textsc{Ecological interpretation}\\
\midrule
$R$          & $10$    & Radius       & Neighborhood radius. Medium range --- neither the slow giants at $R=20$ nor the twitchy sprinters at $R=5$.\\
$T$          & $10$    & Time         & $dt=1/T=0.1$. Moderate temporal sampling; the breath cycle is visible at real time.\\
$\beta_{1}$  & $1/12$  & Inner ring   & Weight $0.083$. Barely a whisper --- this ring is almost acoustically silent.\\
$\beta_{2}$  & $1/6$   & Middle ring  & Weight $0.167$. Dim --- secondary structural sensing.\\
\rowcolor{codebg}$\boldsymbol{\beta_{3}}$ & $\mathbf{1}$ & \textbf{Outer ring} & \textbf{Dominant.} $12\times$ inner, $6\times$ middle. \textit{The organism \emph{is} its surface.}\\
$\mu$        & $0.18$  & Growth centre & Each cell ``wants'' average activation $0.18$ in its outer-ring neighborhood to be in homeostasis.\\
$\sigma$     & $0.033$ & Tolerance    & \textsc{Very} narrow: $G=0$ at $u=0.147$ and $u=0.213$. Window $=0.066$ units. Among the tightest in the taxonomy.\\
\bottomrule
\end{tabularx}
\caption{\small Canonical parameters for \DV\ from \citet{chan2020},
\texttt{animals4D.json}, code 3\textsc{Hy}2v. The highlighted row
$\beta_{3}=1$ is the defining constraint of the species.}
\end{table}

\subsection{The shell kernel formula}

The kernel value at normalized radius $r\in[0,1)$ from the organism's
centre is a three-shell partition of Chan's standard exponential bump:
\begin{equation}
K(r) \;=\; \beta_{k}\cdot
\exp\!\Bigl(\,4-\frac{4}{4\,t\,(1-t)}\Bigr),
\qquad
\begin{cases}
k = \lfloor r n \rfloor & \text{(ring index)}\\
t = r n - k & \text{(local position in ring)}\\
n = 3 & \text{(number of shells)}
\end{cases}
\end{equation}
with $\boldsymbol{\beta}=(\tfrac{1}{12},\tfrac{1}{6},1)$ and outer-ring
dominance $\beta_{3}/\beta_{1}=12$. The core function $\exp\!\bigl(4 -
4/(4r(1-r))\bigr)$ is zero at $r=0$ and $r=1$ (ring boundaries), peaks at
$r=0.5$ within each partition, and is infinitely smooth everywhere. After
normalization (dividing by the sum of all kernel weights over the grid),
the outer ring's \emph{dominance} is preserved but the kernel itself
integrates to unity over the full neighborhood, so the potential field
$U$ remains in $[0,1]$ regardless of the $\beta$ distribution.

\subsection{Acoustic silence of the interior}

A cell at the centre of the organism ($r=0$ relative to its neighbors'
positions) contributes almost nothing to those neighbors' potential
computations --- because the kernel value at $r\approx 0$ (inner ring,
weight $1/12$) is tiny. Conversely, a cell at the organism's surface
($r\approx 5R/6\approx 8.3$ cells from any given internal cell)
contributes enormously, with kernel weight near $1.0$.

This is the defining biology of \DV: \textbf{the skin does all the
sensing.} The interior is nearly acoustically isolated from its own
neighborhood potential. A useful analogy: imagine a starfish whose neural
network lives entirely in its arm tips and whose body core can barely feel
itself. \DV\ experiences its own existence at its boundary, not its
centre. The hyperspherical shell geometry enforces this --- in 4D, surface
area scales as $r^{3}$ while volume scales as $r^{4}$, making the surface
an even larger fraction of the total structure than it is in 3D.

\subsection{The narrow survival band}

The growth function for Lenia is the standard Gaussian-difference form
\begin{equation}
G(u) \;=\; 2\,\exp\!\Bigl(-\frac{(u-\mu)^{2}}{2\sigma^{2}}\Bigr) - 1.
\end{equation}
For \DV, $\mu=0.18$, $\sigma=0.033$, yielding zero-crossings
\begin{equation}
u_{\text{low,\,high}} \;=\; \mu \pm \sigma\sqrt{2\ln 2}
\;=\; 0.147,\ 0.213,
\end{equation}
i.e.\ a survival window of $0.066$ units in potential space. This is
narrow but not extreme among named species: \spc{Orbium} at $\sigma=0.017$
has a window of $\sim 0.023$; \spc{Ignis} at $\sigma=0.012$, a window of
$\sim 0.016$ (among the tightest on record).

What distinguishes \DV's situation is that this narrow band must be
maintained simultaneously across \emph{all 2D cross-sections of a 4D body}.
In a 4D grid, a hyperspherical pattern generates potential gradients in
four independent directions at once. The pattern can only sustain itself
if its 4D geometry keeps every cell's potential within $[0.147, 0.213]$
at all times despite the four-dimensional interference of its own shell
rings. This is the constraint that makes \DV\ a 4D organism rather than
a 2D one.

\section{Four-Dimensional Biology}

A 4D Lenia organism exists on a toroidal grid of four spatial dimensions.
Its body is a \emph{soliton} --- a self-sustaining spatially localized
pattern --- in four-dimensional continuous space. Observation requires
choosing a cutting hyperplane: a 2D subspace of the 4D space. The
cross-section depends entirely on which hyperplane is chosen and at what
orientation.

\subsection{Hyperspherical shell structure}

\DV's 4D body is nested 3-spheres ($S^{3}$): the same relationship that
holds between a circle ($S^{1}$) and a disk in 2D, or a sphere ($S^{2}$)
and a ball in 3D. The three $\beta$ rings correspond to three nested
$S^{3}$ surfaces:

\begin{table}[h!]
\small\sffamily
\renewcommand{\arraystretch}{1.25}
\begin{tabularx}{\linewidth}{@{}l c c >{\raggedright\arraybackslash}X l@{}}
\toprule
\textsc{Shell} & \textsc{4D Radius} & \textsc{$\beta$ Weight} &
\textsc{2D Cross-section} & \textsc{Role}\\
\midrule
\rowcolor{codebg}
Outer $S^{3}$  & $\sim 5R/6$ & $1.000$ & Bright ring at equator      & Dominant sensing surface\\
Middle $S^{3}$ & $\sim R/2$  & $0.167$ & Faint inner ring            & Secondary structure\\
Inner $S^{3}$  & $\sim R/6$  & $0.083$ & Barely visible disk         & Ghost of a core\\
\bottomrule
\end{tabularx}
\caption{\small 4D shell structure of \DV. The outer $S^{3}$ row is
highlighted because $\beta_{3}=1.0$ makes it the dominant feature --- twelvefold brighter than the inner shell.}
\end{table}

\subsection{The ventilating oscillation --- mechanism}

The \ventilans\ behavior is a limit cycle in the 4D state space. The
mechanism: the hyperspherical body expands slightly, pushing cells at its
outer $S^{3}$ shell into regions of lower potential (the neighborhood
density thins as the shell enlarges). When $U$ drops below $\mu-\sigma$
at those cells, $G$ turns negative; the shell contracts. Overshoot is
guaranteed by the finite time step $dt=0.1$. The system oscillates around
the fixed point where every cell's potential exactly equals $\mu=0.18$.

The period of oscillation depends on the local slope of $G$ at $\mu$ ---
which is $0$ at the centre ($G$ is maximally flat at its peak) --- and on
$dt$. For \DV's parameters the breath cycle completes approximately every
$30$--$40$ simulation steps at $\text{spf}=1$, $dt=0.1$, corresponding to
real time on the order of seconds in the Genesis display running at
$60\,$fps.

\subsection{Dimensional comparison}

\begin{table}[h!]
\small\sffamily
\renewcommand{\arraystretch}{1.2}
\begin{tabularx}{\linewidth}{@{}l c c c@{}}
\toprule
\textsc{Property} & \textit{D.\ ventilans} (4D) & \textit{Hexalapillium} (2D) & Typical \textit{Orbium} (2D)\\
\midrule
Dimension      & \cellcolor{codebg}4D toroidal & 2D toroidal    & 2D toroidal\\
Behavior       & \cellcolor{codebg}\SO: \ventilans & \SO: \ventilans & \textsc{to}: translating\\
Symmetry       & \cellcolor{codebg}$O(4)$ hyper & Radial         & Bilateral\\
Kernel peaks   & \cellcolor{codebg}3 ($\beta$-vector) & 1 (single) & 1 (single)\\
$\beta_{3}$    & \cellcolor{codebg}$1.0$ (dominant) & $1.0$       & $1.0$\\
$\sigma$       & \cellcolor{codebg}$0.033$ & $\sim\!0.020$       & $0.017$\\
$R$            & \cellcolor{codebg}$10$   & $\sim\!8$           & $13$\\
Cross-section  & \cellcolor{codebg}Ring $\to$ point & Disk       & Disk\\
Rarity         & \cellcolor{codebg}Extremely rare & Rare         & Common\\
\bottomrule
\end{tabularx}
\caption{\small Comparison of \DV\ with related \SO\ species. Shaded column:
key distinction --- $O(4)$ symmetry makes \DV\ orientation-free. No
preferred axis in 4D.}
\end{table}

\section{Simulation Architecture}

A true 4D Lenia simulation on a $256^{4}$ grid would require
$4{,}294{,}967{,}296$ cells --- computationally intractable in real-time
WebGL2. Genesis uses an \emph{analytical 4D hyperrotation} approach:
\DV's density is approximated analytically from its known shell geometry,
and its temporal evolution is modeled as rigid 4D rotation. This
correctly produces the ventilating oscillation in 2D projection.

\subsection{HYPER\_FRAG shader pipeline}

\begin{lstlisting}[language=glsl, caption={\textsc{hyper\_frag} --- WebGL2 fragment shader, runs per pixel per frame.}]
// Step 1: Map UV pixel to centered 4D position
vec2 xy = (v_uv - 0.5) * 2.2;           // scale to [-1.1, 1.1]
vec4 p4 = vec4(xy.x, xy.y, 0.0, u_wSlice);

// Step 2: Apply 4D rotations (3 independent planes)
p4 = rotXW(p4, u_rotXW);                // lateral tumble X
p4 = rotYW(p4, u_rotYW);                // lateral tumble Y
p4 = rotZW(p4, u_rotZW);                // THE BREATH - ZW rotation

// Step 3: Sample analytical DV density at rotated 4D position
float r     = length(p4) / u_R4D;
float outer = exp(-pow((r - 0.85) / 0.08, 2.0)) * 1.000;
float mid   = exp(-pow((r - 0.55) / 0.10, 2.0)) * 0.167;
float inner = exp(-pow((r - 0.20) / 0.12, 2.0)) * 0.083;
float density = clamp(outer + mid + inner, 0.0, 1.0);

outColor = vec4(density, density, density, 1.0);  // ch3 texture
\end{lstlisting}

The \textsc{zw} rotation \emph{is} the breathing. As \texttt{u\_rotZW}
accumulates, the 4D body's $W$-axis sweeps through the observation plane
and the cross-section ring pulses in and out exactly as the true 4D \SO\
dynamics would produce. The \textsc{xw} and \textsc{yw} rotations are
smaller and add the lateral tumbling characteristic of full 4D dynamics.
Setting all rotation rates to zero freezes \DV\ at a static cross-section,
showing only its equatorial ring structure.

\subsection{The 4D rotation matrices}

\begin{lstlisting}[language=glsl, caption={4D rotation matrices in \textsc{glsl}. The \textsc{zw} plane carries the breath; \textsc{xw}/\textsc{yw} add asymmetric tumble.}]
// Rotation in XW plane (x-axis tilts into W dimension)
vec4 rotXW(vec4 p, float a) {
  return vec4(cos(a)*p.x - sin(a)*p.w,
              p.y, p.z,
              sin(a)*p.x + cos(a)*p.w);
}

// Rotation in ZW plane - primary breathing rotation
vec4 rotZW(vec4 p, float a) {
  return vec4(p.x, p.y,
              cos(a)*p.z - sin(a)*p.w,
              sin(a)*p.z + cos(a)*p.w);
}

// Accumulated angle each frame: theta += speed * 0.016
// Default ZW speed = 0.18 rad/s  =>  ~35s per full rotation
\end{lstlisting}

\subsection{Interface controls}

\begin{table}[h!]
\small\sffamily
\renewcommand{\arraystretch}{1.25}
\begin{tabularx}{\linewidth}{@{}l l l >{\raggedright\arraybackslash}X@{}}
\toprule
\textsc{Control} & \textsc{Range} & \textsc{Default} & \textsc{Effect}\\
\midrule
\textsc{zw} rotation & $0$–$0.8$   & $0.18$ rad/s & Breathing rate. $0.18 \approx 1$ breath / $35$ sim-sec\\
\textsc{xw} rotation & $0$–$0.4$   & $0.05$ rad/s & Lateral $X$ tumble --- asymmetric cross-sections\\
\textsc{yw} rotation & $0$–$0.4$   & $0.07$ rad/s & Lateral $Y$ tumble --- breaks circular symmetry\\
$W$ slice            & $-1.0$–$1.0$ & $0.0$        & Resting hyperplane. $0$ = equatorial (max visibility)\\
4D amplitude         & $0$–$1.5$   & $0.65$        & Density of projected 4D shadow in ch3\\
4D bleed             & $0$–$0.5$   & $0.12$        & Coupling from ch3 into prey ch0 --- seeding strength\\
\bottomrule
\end{tabularx}
\caption{\small 4D simulation interface controls in Genesis. The 4D bleed
parameter is the causal bridge from the 4D body to the 2D ecosystem.}
\end{table}

\section{Multi-Channel Ecosystem Integration}

The Genesis Lenia Expanded Universe substrate runs four coupled channels
simultaneously on a shared $256\times 256$ RGBA float texture (one
\texttt{RGBA32F} pixel per simulation cell):

\begin{table}[h!]
\small\sffamily
\renewcommand{\arraystretch}{1.25}
\begin{tabularx}{\linewidth}{@{}l l l >{\raggedright\arraybackslash}X@{}}
\toprule
\textsc{Channel} & \textsc{Color} & \textsc{Species} & \textsc{Kernel \& parameters}\\
\midrule
R (ch0)    & Gold / amber       & Prey (\spc{Orbium}-class) & $R{=}13$,\ $\beta{=}(1)$,\ $\mu{=}0.15$,\ $\sigma{=}0.017$\\
G (ch1)    & Electric cyan      & Predator (\spc{Ignis}-like) & $R{=}15$,\ $\beta{=}(\tfrac{1}{3},\tfrac{2}{3},1)$,\ $\mu{=}0.26$,\ $\sigma{=}0.036$\\
B (ch2)    & Deep teal          & Morphogen (diffuse)     & $R{=}20$,\ $\beta{=}(1,0.5,0.1)$,\ $\mu{=}0.15$,\ $\sigma{=}0.028$\\
\rowcolor{codebg}
A (ch3)    & Violet / white     & \textit{D.\ ventilans} (4D) & Analytical hyperrotation, $R_{\!4\mathrm{D}}{=}0.85$, $\omega_{\!\mathrm{zw}}{=}0.18$\\
\bottomrule
\end{tabularx}
\caption{\small Four-channel state layout. \DV\ is the only channel
computed analytically rather than via convolution.}
\end{table}

\subsection{Cross-channel coupling equations}

\begin{lstlisting}[language=pseudo, caption={\textsc{sim\_frag} --- the full update rule, per simulation step. \texttt{hyperMix} is the causal dial from 4D shadow to 2D ecology.}]
// Channel 0: Prey
U0 = K0 * A0                                // convolve prey kernel
G0 = grow(U0, mu0, sigma0_eff)              // growth at effective sigma
G0 -= c01 * A1                              // predator suppresses prey
G0 += A3 * hyperMix * 0.4                   // 4D shadow seeds nucleation

// Channel 1: Predator
U1 = K1 * A1                                // convolve predator kernel
G1 = grow(U1, mu1, sigma1_eff)
G1 += c10 * A0 - 0.012                      // prey feeds predator; starvation

// Channel 2: Morphogen
U2 = K2 * A2                                // wide diffuse convolution
sigma0_eff = sigma0 * (1 + c20 * (A2 - 0.3))// morphogen widens prey niche

// Advection (applied before convolution, all channels)
phi = FLOW_FRAG(A0, A2)                     // velocity from gradients
A   = advect(A, grad(A) dot phi)            // Eulerian backward advection
\end{lstlisting}

The \texttt{hyperMix} parameter is the causal bridge between the 4D and 2D
worlds. At $\texttt{hyperMix}=0$, ch3 (\DV) is cosmetic --- a violet
overlay with no coupling. As \texttt{hyperMix} increases, the 4D shadow
actively nudges prey activation wherever the rotating \DV\ body intersects
the 2D plane. Since this intersection pulses spatially and temporally as
\textsc{zw} rotation proceeds, the prey nucleation sites move around the
field in synchrony with \DV's breath.

\subsection{The DV Seed preset --- emergence from hyperspace}

The most striking demonstration is the \textsc{dv seed} initial condition:
all 2D channels begin empty. Only ch3 is initialized with the analytical
\DV\ ring structure at the grid centre.

\begin{lstlisting}[language=pseudo, caption={\textsc{dv seed} initial condition (\texttt{buildEcosystem('hyperseed')}). Only ch3 (alpha) receives non-zero values; the rest of the ecosystem bootstraps from the $\beta_{3}$ ring.}]
for each pixel (x, y) in [0, N)^2:
  nx = (x / N - 0.5) * 2.2                  // centered, normalized
  ny = (y / N - 0.5) * 2.2
  r  = sqrt(nx^2 + ny^2)

  // Three DV rings at their beta-weighted radii
  ring3 = exp(-((r - 0.85) / 0.08)^2) * 1.000
  ring2 = exp(-((r - 0.55) / 0.10)^2) * 0.167
  ring1 = exp(-((r - 0.20) / 0.12)^2) * 0.083

  data[(y*N + x)*4 + 3] = ring3 + ring2 + ring1   // ch3 only
  // ch0 = ch1 = ch2 = 0.0 (empty ecosystem)
\end{lstlisting}

From this --- a single glowing ring of 4D origin, no prey, no predator,
no morphogen --- the ecosystem bootstraps itself. As \textsc{zw} rotation
proceeds, the outer ring ($\beta_{3}=1$, the dominant shell) appears and
disappears as it sweeps through the 2D observation plane. Wherever it
appears, prey nucleates from the \texttt{hyperMix} coupling. Predators do
not appear spontaneously in this preset --- they require separate seeding
or painting. But the prey field that emerges is shaped, from its first
moment of existence, by the geometry of a creature that lives four
dimensions away.

\begin{callout}[Causal transparency \,\textperiodcentered\, \textsc{claire} alignment argument]
Every causal pathway from \DV\ (4D) to the ecosystem (2D) is fully visible
and adjustable in the interface. The \texttt{hyperMix} slider is a causal
dial. The \textit{4D Projection} view shows exactly what \DV\ is doing at
every frame. The \textit{Flow Field} view shows the velocity field
coupling prey and morphogen. There are no hidden variables.

This is the \textsc{claire} principle made literal: Causal, Legible,
Auditable, Interpretable, Robust, Explainable. A 4D organism influences
a 2D world --- and the influence is completely observable.
\end{callout}

\section{Flow Field Advection}

Inspired by Flow-Lenia \citep{plantec2022}, Genesis adds a velocity field
derived from channel gradients that advects the entire state before each
convolution. This adds large-scale fluid dynamics on top of the local
reaction--diffusion.

\begin{lstlisting}[language=pseudo, caption={\textsc{flow\_frag} --- three advection modes. The spiral mode produces galaxy-like large-scale structure.}]
// Mode 0: Gradient flow
// Prey flows along morphogen gradient; prey curl generates vorticity
vel = grad(A2) * 3.0 + curl(grad(A0)) * 1.5

// Mode 1: Curl flow (purely rotational)
vel = curl(grad(A0)) * 4.0

// Mode 2: Spiral flow
angle = atan(centered.y, centered.x) + u_time * 0.3
vel   = radial(centered) * r * 2.0 + grad(A2) * 2.0

// Pack into [0,1] (0.5 = zero velocity), output to flowTex
outColor = vec4(vel * 0.15 + 0.5, 0.5, 1.0)

// Applied in SIM_FRAG as backward advection:
advUV = fract(v_uv - vel * texel * u_flowStr)
prev  = texture(u_state, advUV)
\end{lstlisting}

The flow field makes all channels \emph{fluid}: prey concentrations drift
along morphogen gradients, predator populations are carried in the wake
of prey motion, and the morphogen field itself swirls. In \textit{Spiral}
mode, the entire simulation develops galaxy-like large-scale structure
--- rotating arms of prey with predators chasing along the trailing edges.

\section{Behavioral Observations}

\begin{table}[h!]
\small\sffamily
\renewcommand{\arraystretch}{1.25}
\begin{tabularx}{\linewidth}{@{}l >{\raggedright\arraybackslash}X@{}}
\toprule
\textsc{Preset} & \textsc{Key observation}\\
\midrule
Duel         & Prey solitons orbit then merge under predator pressure. Predator collapses when prey exhausted; reseeding required. Morphogen trails persist $40$–$60$ frames post-death.\\
Swarm        & Arms race: prey cluster in morphogen-depleted voids to evade predator. Predators converge on highest-density prey patches. Stable coexistence occasionally at $c_{01}\!\sim\!0.3$, $c_{10}\!\sim\!0.4$.\\
Coexist      & Separated factions develop distinct morphogen territories. At contact boundary: dramatic species interaction --- sometimes competitive exclusion, occasionally stable interface.\\
Invasion     & Predator wavefront propagates $\sim 2$–$3$ cells / sim-second. Burned region behind front resists recolonization --- morphogen depleted. Prey occasionally outflank via toroidal wrap.\\
\rowcolor{codebg}\textsc{dv} Seed    & Prey nucleates in annular band at outer \DV\ ring radius ($\sim 85\%$ of grid radius). Band rotates with \textsc{zw} speed. \textsc{xw/yw} rotation breaks symmetry, creates asymmetric prey patches. \DV\ is \textit{invisible} in ecosystem view.\\
4D View      & \DV\ ring maximally visible at $W\text{-slice}=0$. Vanishes near $W=\pm 0.8$. Three $\beta$-weighted rings ($R/3$, $R/2$, $5R/6$) clearly distinguishable. Outer ring $\sim 12\times$ brighter than inner --- dominant as expected.\\
Flow: Curl   & Prey generates vorticity; spiral structures self-organize. Predators orbit prey cores. Stable rotating assemblies occasionally form and persist for hundreds of frames.\\
Flow: Spiral & Galaxy-like large-scale structure. Rotating arms of prey, predators trailing along leading edges. \DV\ ring seeds new prey in the arm that sweeps over its position.\\
\bottomrule
\end{tabularx}
\caption{\small Behavioral observations by preset and mode. All observations
at default parameters; coupling strengths $c_{01}=0.35$, $c_{10}=0.40$.}
\end{table}

\section{Open Questions}

\subsection{True 4D simulation}
A full 4D Lenia grid simulation --- impractical at $N=256$ but possible
at $N=32$–$64$ --- would let \DV\ dynamically respond to its own projected
influence on the 2D world. The analytical approximation fixes \DV's shape;
a true simulation might reveal that the 4D--2D coupling warps \DV's own
4D body over time: \emph{an organism shaped by its own shadow.}

\subsection{Phase locking}
If the \DV\ breath period matches the prey--predator Lotka--Volterra
period, resonance effects might produce dramatically amplified prey
bursts at the moment of maximum \DV\ visibility. Can \texttt{hyperMix}
tuning phase-lock these two oscillation systems?

\subsection{Other 4D species}
Chan's \texttt{animals4D.json} contains multiple named 4D organisms. Do
any exhibit translational dynamics in 4D that would produce a moving 2D
nucleation point --- a traveling seed source rather than a stationary
ring? Could such an organism \emph{walk} across the 2D ecosystem, leaving
trails of spontaneous prey?

\subsection{2D relatives}
Does $\boldsymbol{\beta}=(\tfrac{1}{12},\tfrac{1}{6},1)$ produce stable
\SO\ species in 2D Lenia at modified $\mu$–$\sigma$ values? If so, is
there a continuous parameter path from those 2D species to \DV\ --- a
topological deformation across dimensional boundaries?

\subsection{Morphogenetic memory}
At high \texttt{hyperMix} values, does \DV's rotation generate a
\emph{persistent} morphogen ring (via prey secretion) that outlasts the
\DV\ cross-section's visibility? If so, \DV\ leaves a physical record in
a field that modifies the physics of every organism that passes through
it --- a kind of 4D memory inscribed in 2D matter.

\subsection{Synchronization across multiple DV instances}
If two analytical \DV\ generators are seeded with different \textsc{zw}
phase offsets, do the prey populations they nucleate exhibit interference
patterns --- \emph{constructive} (prey burst) when both rings overlap,
\emph{destructive} (prey famine) when they cancel? This would be 4D
organism interference visible in 2D ecology.

\section{Teármann Research Context}

\subsection{The Layer 4 thesis}
The Teármann theoretical framework proposes \textit{Layer 4} reasoning:
the ability to reason \emph{across} causal graph structures rather than
\emph{within} a fixed one --- counterfactual reasoning that operates at
the level of the structural causal model, treating the graph structure
itself as a variable. \citet{pearl2009}'s do-calculus and
\citet{correa2020}'s ctf-calculus operate within a given graph; Layer 4
reasons about which graph to use.

\DV\ is a \emph{physical instantiation} of this concept. The organism
exists in a causal structure (4D Lenia) that is inaccessible to 2D
observers. Its influence on the 2D ecosystem is fully real but cannot be
represented within the 2D causal graph --- \DV\ is a variable that
belongs to a higher-level causal model. Layer 4 reasoning would be
capable of explicitly representing the 4D structure and intervening on
the 4D--2D coupling.

\subsection{Shoal-Broadcast and the morphogen channel}
The Shoal-Broadcast architecture proposes that agents communicate not
through discrete messages but through perturbations in a shared
continuous scalar field --- the \emph{shoal field}. The morphogen channel
(ch2) in Genesis is exactly this: a continuous medium secreted by prey
and predator alike, modulating the physics of survival for every organism
that moves through it. \DV\ adds a higher-order broadcast layer: the 4D
shadow modulates ch0 (prey) which modulates ch2 (morphogen) which
modulates the physics of ch0 and ch1 --- a three-level broadcast cascade
originating in 4D.

\subsection{Rukha and the counterfactual engine}
Rukha (a companion project) is a counterfactual narrative engine grounded
in Pearl's do-calculus and the Correa--Bareinboim ctf-calculus. The \DV\
ecosystem generates exactly the kind of causally transparent data Rukha
is designed to annotate: a sequence of events (prey nucleation, predator
convergence, morphogen trail formation) with completely observable causal
provenance. Each frame is a data point for a causal event graph where the
4D seeding variable (\texttt{ch3}$\to$\texttt{ch0}) is fully observable
and quantifiable via the \texttt{hyperMix} value.

\vspace{22pt}
\begin{center}
{\Large\itshape\color{accent}乙超球 \,\textperiodcentered\, \textit{Dihypersphaerome ventilans}}\\[3pt]
{\sffamily\footnotesize\color{muted}
  \SO\ \,\textperiodcentered\, \ventilans\ \,\textperiodcentered\, $R{=}10$ \,\textperiodcentered\,
  $\boldsymbol{\beta}{=}(\tfrac{1}{12},\tfrac{1}{6},1)$ \,\textperiodcentered\, $\mu{=}0.18$ \,\textperiodcentered\, $\sigma{=}0.033$\\
  4D hyperspherical oscillator \,\textperiodcentered\, \citealt{chan2020} \,\textperiodcentered\,
  code 3\textsc{Hy}2v \,\textperiodcentered\, Third Space AI 2026}
\end{center}

\appendix
\section*{\hfil\normalsize\scshape Mathematical Appendix \hfil}
\addcontentsline{toc}{section}{Mathematical Appendix}
\renewcommand{\thesubsection}{\Alph{subsection}}
\setcounter{subsection}{0}

\subsection{Kernel normalization}
The raw kernel weights must be normalized so that $K$ integrates to unity
over the full neighborhood (i.e.\ a uniform state $A\equiv 1$ yields
$U\equiv 1$). For a discrete 2D grid at resolution $\Delta x = 1/R$,

\begin{lstlisting}[language=pseudo, caption={Kernel normalization.}]
sum = 0
for dy in [-R, R]:
  for dx in [-R, R]:
    r = sqrt(dx^2 + dy^2) / R
    if 0 < r < 1:
      ri = floor(r * 3)
      t  = r * 3 - ri
      sum += beta[ri] * kcore(t)

K_normalized(r) = K_raw(r) / sum

// For DV params: sum is dominated by the outer ring (beta_3 = 1).
// Inner rings contribute < 14% of total weight combined.
\end{lstlisting}

\subsection{Growth-function fixed points}
The growth function $G(u)=2\exp\!\bigl(-(u-\mu)^{2}/2\sigma^{2}\bigr)-1$
has three fixed points in $[0,1]$: $G=0$ at $u=\mu\pm
\sigma\sqrt{2\ln 2}$. For \DV: $u_{\text{eq}} = 0.18 \pm 0.033\cdot
1.1774 = \{0.147,\,0.213\}$. A cell is in equilibrium (neither growing
nor dissolving) when its neighborhood potential lies exactly in this
pair. The organism's stable breathing cycle is a limit cycle that
oscillates around the equilibrium manifold $\{U(x) = 0.18 \text{ for
all } x \text{ in organism body}\}$.

\subsection{4D volume and surface scaling}
In $n$ dimensions, for a shell at radius $r$ with thickness $\mathrm{d}r$:
\begin{equation}
\mathrm{d}V_{n} \;=\; C_{n}\, r^{n-1}\,\mathrm{d}r,
\qquad
C_{n} \;=\; \frac{2\pi^{n/2}}{\Gamma(n/2)}.
\end{equation}
For $n=4$: $C_{4}=2\pi^{2}$. The surface area of \DV's outer $S^{3}$
shell is $A_{3}=2\pi^{2}r^{3}$ --- scaling as $r^{3}$ --- while the 4D
ball volume scales as $r^{4}$. At $r=5R/6$ for \DV\ with $R=10$, the
outer shell radius is $\sim 8.33$ grid cells and its ``4D surface area''
(in grid-cell units) is $\sim 3625$ cells, versus the total 4D ball
volume of $\sim 3801$ cells. The organism is almost entirely surface ---
the interior volume is negligible relative to the 4D surface it presents
to the kernel.

% ============================================================
%  REFERENCES
% ============================================================
\begin{thebibliography}{99}
\small

\bibitem[Chan(2019){Chan}]{chan2019}
Chan, B.\,W.-C. (2019).
\newblock Lenia: Biology of Artificial Life.
\newblock \textit{Complex Systems} 28(3), 251--286.

\bibitem[Chan(2020){Chan}]{chan2020}
Chan, B.\,W.-C. (2020).
\newblock Lenia and Expanded Universe.
\newblock In \textit{Proceedings of the 2020 Conference on Artificial Life}
  (ALIFE 2020), 221--229. MIT Press.

\bibitem[Chan(2020a){Chan}]{chan-animals}
Chan, B.\,W.-C.
\newblock \texttt{animals4D.json}. GitHub: \url{github.com/Chakazul/Lenia}.
\newblock Species code 3\textsc{Hy}2v, Chinese designation 乙超球.

\bibitem[Plantec et al.(2022){Plantec et al.}]{plantec2022}
Plantec, E.\ et al.\ (2022).
\newblock Flow-Lenia: Towards open-ended evolution in cellular automata
  through mass conservation and parameter localization.
\newblock \texttt{arXiv:2212.07906}. ALIFE 2023 Best Paper Award.

\bibitem[Michel et al.(2025){Michel et al.}]{michel2025}
Michel, G.\ et al.\ (2025).
\newblock Exploring Flow-Lenia Universes with a Curiosity-driven AI Scientist.
\newblock \texttt{arXiv:2505.15998}.

\bibitem[Hamon et al.(2024){Hamon et al.}]{hamon2024}
Hamon, G.\ et al.\ (2024).
\newblock Discovering self-organized patterns in Lenia with
  curiosity-driven exploration. arXiv preprint.

\bibitem[Faldor et al.(2024){Faldor et al.}]{faldor2024}
Faldor, M.\ et al.\ (2024).
\newblock Toward Artificial Open-Ended Evolution within Lenia using
  Quality-Diversity. \texttt{arXiv:2406.04235}.

\bibitem[Pearl(2009){Pearl}]{pearl2009}
Pearl, J.\ (2009).
\newblock \textit{Causality: Models, Reasoning, and Inference}
  (2nd ed.). Cambridge University Press.

\bibitem[Correa \& Bareinboim(2020){Correa and Bareinboim}]{correa2020}
Correa, J.\,D.\ \& Bareinboim, E.\ (2020).
\newblock A calculus for stochastic interventions: Causal effect
  identification and surrogate experiments.
\newblock In \textit{Proceedings of AAAI 2020}.

\bibitem[Sebastian \& Claude(2026){Sebastian and Claude}]{koons2026}
Sebastian, S.\ \& Claude (2026).
\newblock \textit{Orbium unicaudatus ignis var.\ phantasma}: A Lenia
  Species Engineered to Inhabit the Edge of Chaos. Third Space AI internal
  report.

\bibitem[Rafler(2011){Rafler}]{rafler2011}
Rafler, S.\ (2011).
\newblock Generalization of Conway's Game of Life to a continuous
  domain: SmoothLife. \texttt{arXiv:1111.1567}.

\bibitem[Schmickl et al.(2016){Schmickl et al.}]{schmickl2016}
Schmickl, T.\ et al.\ (2016).
\newblock How a life-like system emerges from a simplistic particle
  motion law. \textit{Scientific Reports} 6.

\bibitem[Quílez(2013){Quílez}]{quilez2013}
Quílez, I.\ (2013).
\newblock Palettes. \url{iquilezles.org/articles/palettes/}.

\end{thebibliography}

\vspace{18pt}
\begin{center}
\sffamily\footnotesize\color{muted}
Third Space AI \,\textperiodcentered\, 2026\\
\href{https://github.com/Kquant03/genesis-phase-transition}{\texttt{github.com/Kquant03/genesis-phase-transition}}
\end{center}

\end{document}
